\chapter{Cryptographic Hash Function: SHA-256}
\label{ch:sha256}

\section{Historical Background, Development, and Design Motivation}
Cryptographic hash functions serve as fundamental building blocks across modern information security, providing data integrity verification, password authentication, and digital signature generation. Early standard hash functions developed in the 1990s, including MD5 (128-bit) and SHA-1 (160-bit), relied on the Merkle-Damg{\aa}rd construction. However, theoretical and practical cryptanalytic breakthroughs severely compromised these algorithms:
\begin{itemize}
    \item In 2004, Xiaoyun Wang et al. demonstrated practical collision attacks against MD5.
    \item In 2017, the SHAttered project (CWI Amsterdam and Google) produced the first real-world collision for full SHA-1 with a computational effort equivalent to $2^{63.1}$ evaluations.
\end{itemize}
Anticipating the theoretical degradation of 160-bit hash designs, the National Security Agency (NSA) designed the SHA-2 family, published by NIST in 2002 under FIPS PUB 180-2 and updated in FIPS PUB 180-4 \cite{fips180_4}. Within this family, \textbf{SHA-256} was specifically engineered to deliver a 128-bit security margin against collision attacks (matching the symmetric strength of AES-128) by employing a 256-bit digest, an expanded 64-round compression function, and non-linear Boolean mixing primitives.

\section{Underlying Mathematical Foundations and Operating Principles}
SHA-256 operates on 32-bit binary words under modular arithmetic in $\mathbb{Z}_{2^{32}}$ combined with bitwise Boolean algebra.

\subsection{Word-Level Operations and Bitwise Primitive Functions}
All variables are 32-bit unsigned integers (range $0 \le X < 2^{32}$). Addition is defined modulo $2^{32}$:
\begin{equation}
X \boxplus Y = (X + Y) \bmod 2^{32}
\end{equation}
The basic bitwise operators include bitwise XOR ($\oplus$), bitwise AND ($\wedge$), bitwise OR ($\vee$), and bitwise NOT ($\neg$ or $\bar{X}$).

The algorithm utilizes right rotation $\text{ROTR}^n(X)$ and right logical shift $\text{SHR}^n(X)$:
\begin{align}
\text{ROTR}^n(X) &= (X \gg n) \vee (X \ll (32 - n)) \\
\text{SHR}^n(X) &= X \gg n
\end{align}

\subsection{Non-Linear Compression Functions}
To guarantee diffusion, confusion, and strict avalanche criteria, SHA-256 employs six non-linear auxiliary functions:
\begin{align}
\text{Ch}(X, Y, Z) &= (X \wedge Y) \oplus (\neg X \wedge Z) \\
\text{Maj}(X, Y, Z) &= (X \wedge Y) \oplus (X \wedge Z) \oplus (Y \wedge Z) \\
\Sigma_0(X) &= \text{ROTR}^2(X) \oplus \text{ROTR}^{13}(X) \oplus \text{ROTR}^{22}(X) \\
\Sigma_1(X) &= \text{ROTR}^6(X) \oplus \text{ROTR}^{11}(X) \oplus \text{ROTR}^{25}(X) \\
\sigma_0(X) &= \text{ROTR}^7(X) \oplus \text{ROTR}^{18}(X) \oplus \text{SHR}^3(X) \\
\sigma_1(X) &= \text{ROTR}^{17}(X) \oplus \text{ROTR}^{19}(X) \oplus \text{SHR}^{10}(X)
\end{align}
The function $\text{Ch}$ (``choose'') selects bits from $Y$ if the corresponding bit of $X$ is 1, and from $Z$ otherwise. The function $\text{Maj}$ (``majority'') outputs 1 if two or three of the input bits are 1.

\subsection{Mathematical Derivation of Constants}
In accordance with the ``nothing up my sleeve'' security principle, the 8 initial hash values $H_0^{(0)} \dots H_7^{(0)}$ and the 64 round constants $K_0 \dots K_{63}$ are derived mathematically:
\begin{itemize}
    \item $H^{(0)}$: The first 32 bits of the fractional parts of the square roots of the first 8 prime numbers ($2, 3, \dots, 19$):
    \begin{equation}
    H_i^{(0)} = \lfloor (\sqrt{p_i} - \lfloor\sqrt{p_i}\rfloor) \times 2^{32} \rfloor
    \end{equation}
    \item $K_t$: The first 32 bits of the fractional parts of the cube roots of the first 64 prime numbers ($2, 3, \dots, 311$):
    \begin{equation}
    K_t = \lfloor (\sqrt[3]{p_t} - \lfloor\sqrt[3]{p_t}\rfloor) \times 2^{32} \rfloor
    \end{equation}
\end{itemize}

\section{Complete Operational Procedure}
The overall execution pipeline of SHA-256 consists of Message Preprocessing (Padding) followed by iterative block compression following the Merkle-Damg{\aa}rd paradigm, as illustrated in Figure \ref{fig:sha_pipeline}.

\begin{figure}[htbp]
\centering
\begin{tikzpicture}[>=stealth,
    block/.style={rectangle, draw=blue!70!black, fill=blue!5, thick, text width=9.2cm, align=center, rounded corners=3pt, inner sep=6pt},
    io/.style={rectangle, draw=black!70, fill=gray!12, thick, text width=7.0cm, align=center, rounded corners=3pt, inner sep=5pt},
    line/.style={draw, -latex, thick}]
    
    \node (msg) [io] {\textbf{Arbitrary-Length Input Message} $M$ ($L$ bits)};
    \node (pad) [block, below=0.45cm of msg] {\textbf{Preprocessing \& Padding Phase:}\\
        \vspace{2pt}
        \small Append \texttt{0x80} $\to$ zero padding $\to$ 64-bit length $L$ ($N \times 512$ bits)};
    \node (comp) [block, below=0.45cm of pad] {\textbf{Davies-Meyer Compression Function (64 Rounds per block):}\\
        \vspace{2pt}
        \small $H^{(i)} = H^{(i-1)} \boxplus \text{Compress}(H^{(i-1)}, M^{(i)})$};
    \node (out) [io, below=0.45cm of comp] {\textbf{Final 256-bit Digest Output:}\\
        \small $H_0^{(N)} \parallel H_1^{(N)} \parallel \dots \parallel H_7^{(N)}$};
        
    \path [line] (msg) -- (pad);
    \path [line] (pad) -- (comp);
    \path [line] (comp) -- (out);
\end{tikzpicture}
\caption{SHA-256 Merkle-Damg{\aa}rd Construction and Iterative Compression Pipeline.}
\label{fig:sha_pipeline}
\end{figure}

\subsection{Message Preprocessing and Padding}
Given an input bit-string $M$ of length $L$ bits:
\begin{enumerate}
    \item Append a single bit `1' (hexadecimal \texttt{0x80}).
    \item Append $k$ zero bits `0', where $k$ is the smallest non-negative integer satisfying:
    \begin{equation}
    L + 1 + k \equiv 448 \pmod{512}
    \end{equation}
    \item Append the 64-bit big-endian binary representation of the original length $L$.
\end{enumerate}
The resulting padded message has a length exactly equal to a multiple of 512 bits (64 bytes), divisible into $N$ blocks: $M^{(1)}, M^{(2)}, \dots, M^{(N)}$.

\subsection{Compression Function}
For each 512-bit block $M^{(i)}$:
\begin{enumerate}
    \item \textbf{Message Schedule Expansion:} Decompose $M^{(i)}$ into sixteen 32-bit words $W_0 \dots W_{15}$. Expand these into 64 words:
    \begin{equation}
    W_t = 
    \begin{cases}
    M_t^{(i)}, & 0 \le t \le 15 \\
    \sigma_1(W_{t-2}) \boxplus W_{t-7} \boxplus \sigma_0(W_{t-15}) \boxplus W_{t-16}, & 16 \le t \le 63
    \end{cases}
    \end{equation}
    \item \textbf{Initialize Working Registers:}
    \begin{equation}
    (a, b, c, d, e, f, g, h) = (H_0^{(i-1)}, H_1^{(i-1)}, \dots, H_7^{(i-1)})
    \end{equation}
    \item \textbf{64-Round Transformation Loop:} For $t = 0, \dots, 63$:
    \begin{align}
    T_1 &= h \boxplus \Sigma_1(e) \boxplus \text{Ch}(e, f, g) \boxplus K_t \boxplus W_t \\
    T_2 &= \Sigma_0(a) \boxplus \text{Maj}(a, b, c) \\
    h &= g, \quad g = f, \quad f = e, \quad e = d \boxplus T_1 \\
    d &= c, \quad c = b, \quad b = a, \quad a = T_1 \boxplus T_2
    \end{align}
    \item \textbf{State Accumulation (Davies-Meyer Feedforward):}
    \begin{equation}
    H_0^{(i)} = H_0^{(i-1)} \boxplus a, \quad \dots, \quad H_7^{(i)} = H_7^{(i-1)} \boxplus h
    \end{equation}
\end{enumerate}
The final 256-bit digest is the concatenation of $H_0^{(N)} \parallel H_1^{(N)} \parallel \dots \parallel H_7^{(N)}$.

\section{Security Assumptions, Properties, and Known Attacks}
\subsection{Cryptographic Security Properties}
A secure cryptographic hash function must satisfy three canonical properties:
\begin{enumerate}
    \item \textbf{Pre-image Resistance (One-Wayness):} Given hash value $y$, it is computationally infeasible to find any $x$ such that $\text{SHA-256}(x) = y$. Work factor: $2^{256}$.
    \item \textbf{Second Pre-image Resistance (Weak Collision Resistance):} Given input $x$, it is infeasible to find $x' \ne x$ such that $\text{SHA-256}(x') = \text{SHA-256}(x)$. Work factor: $2^{256}$.
    \item \textbf{Collision Resistance (Strong Collision Resistance):} It is infeasible to find any distinct pair $(x, x')$ such that $\text{SHA-256}(x) = \text{SHA-256}(x')$. By the Birthday Paradox, the work factor is $2^{256/2} = 2^{128}$.
\end{enumerate}

\subsection{Known Attacks and Mitigation}
\begin{itemize}
    \item \textbf{Length Extension Attack:} Because SHA-256 uses the Merkle-Damg{\aa}rd structure without final truncation, an attacker knowing $\text{SHA-256}(M)$ and $\text{len}(M)$ can compute $\text{SHA-256}(M \parallel \text{pad} \parallel M')$ without knowing $M$. This renders raw $\text{SHA-256}(k \parallel m)$ insecure for Message Authentication Codes. \emph{Mitigation:} Use HMAC-SHA-256 or SHA-512/256 (which truncates the state).
\end{itemize}

\section{Advantages, Limitations, and Modern Applications}
\subsection{Advantages}
SHA-256 provides a robust 128-bit collision resistance margin, deterministic execution, and excellent performance in hardware and modern 64-bit CPUs via Intel SHA Extensions.

\subsection{Limitations}
Vulnerable to length-extension attacks when improperly used as a raw MAC, and relatively slow for ultra-high-throughput streaming compared to newer tree hashes (like BLAKE3).

\subsection{Modern Applications}
SHA-256 is the integrity foundation of modern computing: TLS 1.3 certificate validation, Bitcoin Proof-of-Work mining, Linux package signing, and Git object identifier version 2.

\section{Complete Worked Example Step-by-Step}
We examine the computation of $\text{SHA-256}(\texttt{"abc"})$:
\begin{enumerate}
    \item \textbf{Input:} \texttt{"abc"} corresponds to 3 bytes: \texttt{0x61, 0x62, 0x63}. Length $L = 24$ bits.
    \item \textbf{Padding:}
    \begin{itemize}
        \item Append \texttt{0x80}.
        \item Append 55 zero bytes (\texttt{0x00}).
        \item Append length 24 as 8 bytes: \texttt{00 00 00 00 00 00 00 18}.
        \item Total block size = 64 bytes (512 bits).
    \end{itemize}
    \item \textbf{Initial Words $W_0 \dots W_{15}$:}
    \begin{itemize}
        \item $W_0 = \texttt{0x61626380}$
        \item $W_1 = W_2 = \dots = W_{14} = \texttt{0x00000000}$
        \item $W_{15} = \texttt{0x00000018}$
    \end{itemize}
    \item \textbf{Expanded Word $W_{16}$:}
    \begin{align}
    W_{16} &= \sigma_1(W_{14}) \boxplus W_9 \boxplus \sigma_0(W_1) \boxplus W_0 \nonumber \\
           &= 0 \boxplus 0 \boxplus 0 \boxplus \texttt{0x61626380} = \texttt{0x61626380}
    \end{align}
    \item \textbf{Round 0 Execution:}
    With initial registers $(a, b, c, d, e, f, g, h) = (\texttt{6a09e667}, \dots, \texttt{5be0cd19})$ and $K_0 = \texttt{0x428a2f98}$:
    \begin{align*}
    T_1 &= h \boxplus \Sigma_1(e) \boxplus \text{Ch}(e,f,g) \boxplus K_0 \boxplus W_0 = \texttt{0x54da50e8} \\
    T_2 &= \Sigma_0(a) \boxplus \text{Maj}(a,b,c) = \texttt{0x08909ae5} \\
    a_{\text{new}} &= T_1 \boxplus T_2 = \texttt{0x5d6aebcd} \\
    e_{\text{new}} &= d \boxplus T_1 = \texttt{0xfa2a4622}
    \end{align*}
    After 64 rounds and state addition, the output matches NIST KAT:
    \begin{equation}
    {\small\text{Digest} = {\scriptsize\texttt{ba7816bf8f01cfea414140de5dae2223b00361a396177a9cb410ff61f20015ad}}}
    \end{equation}
\end{enumerate}

\clearpage
\section{Python Implementation from First Principles}
The implementation avoids \texttt{hashlib}, constructing the round constants dynamically from the first 64 prime cube roots and applying 32-bit bitwise masking, as illustrated in Listing \ref{lst:sha256_code}:
\begin{lstlisting}[language=Python, caption={Core compression loop of SHA-256.}, label={lst:sha256_code}]
def sha256_compress(h_state, block):
    w = [0] * 64
    for t in range(16):
        w[t] = int.from_bytes(block[t*4 : (t+1)*4], "big")
    for t in range(16, 64):
        s0 = small_sigma0(w[t - 15])
        s1 = small_sigma1(w[t - 2])
        w[t] = (w[t - 16] + s0 + w[t - 7] + s1) & MASK_32

    a, b, c, d, e, f, g, h = h_state
    for t in range(64):
        t1 = (h + big_sigma1(e) + ch(e, f, g) + ROUND_CONSTANTS_K[t] + w[t]) & MASK_32
        t2 = (big_sigma0(a) + maj(a, b, c)) & MASK_32
        h, g, f = g, f, e
        e = (d + t1) & MASK_32
        d, c, b = c, b, a
        a = (t1 + t2) & MASK_32
    return [(h_state[i] + val) & MASK_32 for i, val in enumerate([a,b,c,d,e,f,g,h])]
\end{lstlisting}

\section{Experimental Evaluation and Correctness Verification}
The implementation was verified against three official NIST FIPS 180-4 test vectors:
\begin{enumerate}
    \item Standard string \texttt{"abc"}: \textbf{PASSED (100\% match)}.
    \item Empty string \texttt{""}: \textbf{PASSED (100\% match)}.
    \item 448-bit string (two-block padding): \textbf{PASSED (100\% match)}.
\end{enumerate}

\begin{table}[H]
\centering
\caption{SHA-256 Experimental Throughput Benchmark.}
\label{tab:sha_perf}
\begin{tabular}{lcc}
\toprule
\textbf{Buffer Size} & \textbf{Execution Latency} & \textbf{Calculated Throughput} \\
\midrule
1 KB  & 15.76 ms & 63.45 KB/s \\
10 KB & 105.56 ms & 94.73 KB/s \\
50 KB & 414.75 ms & 120.55 KB/s \\
\bottomrule
\end{tabular}
\end{table}
The experimental results demonstrate linear time complexity $\mathcal{O}(N)$ in terms of the number of 512-bit message blocks.

\needspace{19\baselineskip}
\subsection*{AI Assistance and Prompt Provenance}
\begin{itemize}
    \item \textbf{AI-Assisted Scope:} Theory draft, compression logic, worked example, Listing~\ref{lst:sha256_code}, and Table~\ref{tab:sha_perf} (paired with \texttt{src/sha256.py}).
    \item \textbf{Session Steps:} Baseline in \textbf{Step 1} (\texttt{P01}--\texttt{P05}); formatting, Round 0 arithmetic, and benchmarks verified in \textbf{Steps 3, 6, 9}.
    \item \textbf{Student Verification:} Verified all three NIST FIPS 180-4 KAT vectors and throughput measurements against \texttt{evidence/\allowbreak terminal\_output.txt}.
\end{itemize}
\begin{lstlisting}[style=promptstyle, caption={Compact Prompt (R-REP-03) for Chapter 3: SHA-256.}, label={lst:prompt_ch3}]
Following the project specification (Student: TRAN HUU QUYET, SID: 066206014296), write the SHA-256 chapter in English LaTeX covering all 8 mandatory sections, including a step-by-step worked example, pure Python snippets from first principles (without hashlib), and NIST FIPS 180-4 KAT verification and benchmarks.
\end{lstlisting}
